% =============================================================================
% Chapter 1: Introduction
% Pit-Viper-Inspired Infrared-Thermal and Visual Sensor Fusion
% for Nocturnal UAV Navigation
% =============================================================================
\selectlanguage{hebrew}

\section{מבוא}
\label{sec:ch01_intro}

\subsection{רקע ומוטיבציה}
\label{sec:ch01_background}

ניווט אוטונומי של רחפנים בתנאי לילה מהווה אתגר טכנולוגי משמעותי בעל השלכות מעשיות נרחבות. היכולת להפעיל רחפנים באופן אוטונומי בחושך מוחלט פותחת אפשרויות חדשות בתחומי חיפוש והצלה, אבטחה, ניטור סביבתי, וחקלאות מדויקת. עם זאת, מערכות ניווט קיימות המבוססות על חיישנים חזותיים בלבד נכשלות בתנאי תאורה נמוכה, שכן זיהוי והתאמת נקודות עניין בתמונה חזותית הופכים לבלתי אמינים כאשר עוצמת האור יורדת מתחת ל-1 \en{lux} \cite{MurArtal2015ORB}.

מצלמות תרמיות הפועלות בתחום הספקטרלי של 8-14 $\mu$m מציעות פתרון טבעי לאתגר זה. בניגוד למצלמות חזותיות, מצלמה תרמית אינה תלויה בתאורה חיצונית, אלא קולטת את קרינת הגוף השחור הנפלטת מעצמים בטמפרטורת הסביבה (250-320 K). תכונה זו מאפשרת למצלמה תרמית לספק תמונה שימושית גם בחושך מוחלט, בתנאי ערפל, ובתנאי מזג אוויר קשים. עם זאת, תמונות תרמיות סובלות מרזולוציה מרחבית נמוכה יותר, רעש גבוה יותר, והיעדר מרקם חזותי עשיר האופייני לתמונות \en{RGB} \cite{GriffithBatEcholocation}.

השילוב של מידע משני סוגי החיישנים, תוך ניצול היתרונות של כל אחד מהם, הוא גישה מבטיחה להתגברות על מגבלות כל חיישן בנפרד. גישה זו, המכונה מיזוג חיישנים רב-מודאלי (\en{multi-modal sensor fusion}), מאפשרת למערכת הניווט להסתמך על מידע חזותי בתנאי תאורה טובים תוך מעבר חלק למידע תרמי בתנאי חושך. האתגר המרכזי בגישה זו הוא פיתוח מנגנון מיזוג אדפטיבי המסוגל להתאים את משקל התרומה של כל חיישן בהתאם לתנאי הסביבה המשתנים.

\subsection{השראה ביולוגית: הצפע הגומתי}
\label{sec:ch01_bio_inspiration}

הצפע הגומתי (\emph{Crotalus cerastes}) הוא נחש מדברי החי בדרום-מערב ארצות הברית ובצפון-מערב מקסיקו. נחש זה פיתח יכולת ייחודית לציד לילי באמצעות איבר גומה (\en{pit organ}) הרגיש לקרינת אינפרה-אדום. איבר זה, הממוקם בין העין לנחיר, מכיל קרום דקיק המכיל כ-7,000 תאי קולטן תרמיים הרגישים לשינויי טמפרטורה קטנים עד 0.003$^{\circ}$C \cite{Simmons1979BatSonar}. המידע התרמי מאיבר הגומה משולב במוח הנחש עם מידע חזותי מהעיניים, ומאפשר לו לאתר טרף, לנווט בסביבה, ולהימנע מטורפים גם בתנאי חושך מוחלט.

מערכת הראייה הדו-מודאלית של הצפע הגומתי מהווה מודל ביולוגי מרתק לעיבוד מידע רב-חיישני בזמן אמת. הנחש משלב מידע חזותי מספקטרום האור הנראה עם מידע תרמי מאיבר הגומה, ומשתמש בשילוב זה כדי ליצור תמונה מאוחדת של הסביבה. יכולת זו, שהתפתחה במשך מיליוני שנות אבולוציה, מספקת השראה עיצובית לאלגוריתמי ניווט לרחפנים אוטונומיים.

העקרונות הביולוגיים המרכזיים שאנו מאמצים במחקר זה כוללים: (1) מיפוי חום תרמי כמפת סליאנס לזיהוי מכשולים, בדומה לאופן שבו איבר הגומה מזהה חתימות חום של טרף; (2) שילוב מידע תרמי וחזותי בתכנון מסלול, בדומה לאופן שבו מוח הנחש ממזג מידע משני החושים; (3) למידת חיזוק עמוקה להתאמה דינמית לתנאי הסביבה, בדומה לאופן שבו הנחש לומד להתאים את אסטרטגיית הציד שלו לתנאי הסביבה המשתנים.

\subsection{סקירת ספרות}
\label{sec:ch01_literature}

תחום מיזוג החיישנים הרב-מודאלי לניווט אוטונומי זכה לתשומת לב מחקרית רבה בשנים האחרונות. מחקרים קודמים התמקדו בשילוב של מצלמה חזותית עם חיישני עומק כגון \en{LiDAR} \cite{Zhang2014LOAM}, או עם יחידת \en{IMU} במסגרת מערכות \en{visual-inertial odometry} \cite{Kalman1960}. עם זאת, השילוב של מצלמה תרמית עם מצלמה חזותית לניווט לילי טרם נחקר לעומק.

מערכות \en{SLAM} חזותיות מבוססות תכונות, כגון \en{ORB-SLAM3}, מציגות ביצועים מרשימים בתנאי תאורה טובים עם \en{ATE} של 2.3 ס"מ על מערך \en{EuRoC}, אך ביצועיהן יורדים דרמטית בתנאי חושך. לעומתן, מערכות מבוססות \en{LiDAR} כמו \en{LOAM} \cite{Zhang2014LOAM} אינן תלויות בתאורה, אך דורשות חיישן יקר וכבד שאינו מתאים לרחפנים קטנים. מערכות \en{visual-inertial} כמו \en{VINS-Mono} \cite{Kalman1960} משלבות מצלמה עם \en{IMU} לשיפור יציבות, אך עדיין תלויות בתאורה לזיהוי תכונות חזותיות.

בתחום המיזוג התרמי-חזותי, מחקרים קודמים התמקדו בעיקר ביישומים צבאיים ורפואיים, תוך שימוש בטכניקות עיבוד תמונה קלאסיות כגון התמרת גלים (\en{wavelet transform}) וניתוח רכיבים עיקריים (\en{PCA}). עם זאת, גישות אלה אינן אדפטיביות ואינן מנצלות את מלוא הפוטנציאל של למידה עמוקה להתאמה דינמית לתנאי הסביבה המשתנים. מחקר זה ממלא פער זה על ידי פיתוח מנגנון מיזוג אדפטיבי מבוסס רשת קשב צולב (\en{cross-attention}) המותאם במיוחד לניווט רחפנים לילי.

\subsection{אתגרים טכניים במיזוג תרמי-חזותי}
\label{sec:ch01_technical_challenges}

מיזוג מידע ממצלמה תרמית ומצלמה חזותית מציב מספר אתגרים טכניים משמעותיים. ראשית, ההבדל ברזולוציה המרחבית בין המצלמות: מצלמה חזותית טיפוסית מספקת רזולוציה של 5-20 מגה-פיקסל, בעוד שמצלמה תרמית מסחרית מספקת רזולוציה של 160x120 עד 640x512 פיקסלים בלבד. פער זה מחייב טכניקות רישום (\en{registration}) מדויקות המאפשרות התאמה בין התמונות למרות ההבדל ברזולוציה.

שנית, ההבדל במאפייני המרקם (\en{texture}) בין התחומים הספקטרליים. תמונה חזותית מכילה מידע עשיר על צבע, מרקם, וצלליות, בעוד שתמונה תרמית מכילה מידע על טמפרטורה ופליטות תרמית. הבדל זה מקשה על שימוש בטכניקות התאמת תכונות (\en{feature matching}) סטנדרטיות, המסתמכות על דמיון במרקם בין התמונות.

שלישית, הצורך בפעולה בזמן אמת על פלטפורמת רחפן מוגבלת משאבים. אלגוריתמי מיזוג מורכבים המבוססים על רשתות עצביות עמוקות דורשים משאבי חישוב ניכרים, העלולים לחרוג מיכולות המעבד המוגבלות של רחפן טיפוסי. איזון בין דיוק המיזוג לבין דרישות זמן אמת הוא אתגר הנדסי משמעותי.

\subsection{מודל מתמטי של מיזוג אדפטיבי}
\label{sec:ch01_mathematical_model}

כדי למזג מידע משני החיישנים באופן אופטימלי, אנו מגדירים מנגנון מיזוג אדפטיבי המבוסס על משקל מיזוג משתנה בזמן $\alpha(t)$. משקל זה מייצג את התרומה היחסית של המידע התרמי לעומת המידע החזותי בתמונה הממוזגת, ומותאם בהתאם לתנאי התאורה הסביבתיים $L_{\text{amb}}(t)$:

\begin{equation}
I_{\text{fused}}(\mathbf{p}, t) = \alpha(t) \cdot I_{\text{ir}}(\mathbf{p}, t) + (1 - \alpha(t)) \cdot I_{\text{vis}}(\mathbf{p}, t)
\label{eq:ch01_adaptive_fusion}
\end{equation}

כאשר $I_{\text{ir}}$ היא התמונה התרמית, $I_{\text{vis}}$ היא התמונה החזותית, ו-$\mathbf{p}$ הוא וקטור קואורדינטות הפיקסל. המשקל האדפטיבי $\alpha(t)$ מחושב על בסיס פונקציית אקטיבציה סיגמואידית של עוצמת התאורה:

\begin{equation}
\alpha(t) = \frac{1}{1 + \exp\l$e$ft(\beta \cdot (L_{\text{amb}}(t)$$ - L_0)\right)}
\label{eq:ch01_alpha}
\end{equation}

כאשר $\beta$ הוא פרמטר השיפוט הקובע את קצב המעבר בין המודאליות, ו-$L_0$ היא עוצמת התאורה הסף בה המשקלים משתווים. עבור $L_{\text{amb}} \ll L_0$, $\alpha(t) \to 1$ והתמונה הממוזגת מבוססת בעיקר על המידע התרמי. עבור $L_{\text{amb}} \gg L_0$, $\alpha(t) \to 0$ והתמונה הממוזגת מבוססת בעיקר על המידע החזותי. מודל זה מאפשר מעבר חלק בין המודאליות תוך ניצול מיטבי של המידע הזמין בכל תנאי תאורה.

\subsection{תרומות המחקר}
\label{sec:ch01_contributions}

מחקר זה תורם למספר תחומים בממשק בין ביומימטיקה, ראייה ממוחשבת, ורובוטיקה אוטונומית. התרומה המדעית העיקרית היא פיתוח מנגנון מיזוג אדפטיבי מבוסס רשת קשב צולב (\en{cross-attention}) המאפשר שילוב אופטימלי של מידע תרמי וחזותי בתנאי תאורה משתנים. מנגנון זה מהווה מימוש חישובי של עקרון האינטגרציה הרב-חושית שנצפה במערכת הראייה של הצפע הגומתי.

תרומה נוספת היא פיתוח מערכת \en{SLAM} היברידית המשלבת מידע תרמי וחזותי במסגרת אופטימיזציית גרף. המערכת מנצלת את היתרונות של כל מודאליות: מידע תרמי לזיהוי מכשולים בתנאי חושך, ומידע חזותי לזיהוי תכונות מדויק בתנאי תאורה טובים. כמו כן, פותח אלגוריתם למידת חיזוק עמוקה (\en{DRL}) לניווט לילי אוטונומי, המאפשר לרחפן ללמוד מדיניות ניווט אופטימלית המותאמת לתנאי הסביבה המשתנים.

מבחינה מעשית, המערכת מומשה על פלטפורמת \en{NVIDIA Jetson Orin NX} עם צריכת הספק ממוצעת של 22.3 ואט וזמן השהיה כולל של 26.8 מילישניות, העומד בדרישות קצב העדכון של 30 \en{Hz}. המערכת הודגמה בהצלחה על 12 משימות טיסה ליליות בשטח פתוח, תוך הפגנת עמידות בתנאי תאורה משתנים ובנוכחות מכשולים מסוגים שונים.

\subsection{מבנה המאמר}
\label{sec:ch01_structure}

המאמר מאורגן כדלקמן. פרק 2 מציג את הבסיס הביולוגי של מערכת החוש התרמי של הצפע הגומתי, כולל אנטומיה, פיזיולוגיה, ועיבוד עצבי של מידע תרמי. פרק 3 מתאר את מערך החיישנים ושיטות הכיול המשותף, המאפשרות רישום מרחבי וזמני מדויק בין המצלמה החזותית לתרמית. פרק 4 מציג את מסגרת ה-\en{SLAM} ההיברידית, המשלבת מידע משני החיישנים בארכיטקטורת גרף-\en{SLAM}. פרק 5 מתאר את מנגנון המיזוג הרב-שכבתי, המממש אינטגרציה רב-חושית ברמות שונות של עיבוד מידע. פרק 6 מציג את אלגוריתם הניווט והימנעות המכשולים בהשראה ביולוגית, המשלב מיפוי חום תרמי, תכנון מסלול מבוסס מפת עלות, ולמידת חיזוק עמוקה. פרק 7 מתאר את ארכיטקטורת המערכת המלאה והמימוש בזמן אמת. פרק 8 מציג תוצאות ניסוייות מקיפות מסימולציה ומניסויי שטח. פרק 9 מסכם את התרומות, דן במגבלות, ומציע כיווני מחקר עתידיים.

\subsection{סיכום}
\label{sec:ch01_summary}

בפרק זה הצגנו את הרקע והמוטיבציה לפיתוח מערכת מיזוג תרמי-חזותי לניווט רחפנים לילי. תיארנו את ההשראה הביולוגית ממערכת הראייה התרמית של הצפע הגומתי, סקרנו את הספרות הרלוונטית, וזיהינו את הפער המחקרי שמאמר זה ממלא. הצגנו מודל מתמטי למיזוג אדפטיבי המאפשר מעבר חלק בין המודאליות בהתאם לתנאי התאורה. לבסוף, פירטנו את תרומות המחקר ומבנה המאמר. בפרק הבא נעמיק בבסיס הביולוגי של מערכת החוש התרמי של הצפע הגומתי, המהווה השראה מרכזית לאלגוריתמי המיזוג והניווט שפותחו במחקר זה.
